\section{Future Work}
\label{sec:future_work}

Several extensions would meaningfully strengthen the project beyond the current scope. These are ordered by expected return on effort, not by ambition.

\begin{enumerate}
  \item \textbf{Scale the per-series classical and deep studies to the full panel.} The four representative series in \cref{sec:eda:repseries} were chosen to span the panel's stress dimensions, but the 36{,}923-series panel is much wider. Repeating the classical and deep protocols on every eligible series (the 1{,}930 series that cleanly cross the validation cutoff) would produce a per-series leaderboard with proper coverage of the panel rather than a stress-test sample.
  \item \textbf{Per-horizon calibration of the LightGBM quantile method.} \Cref{sec:quantile} fits a single LightGBM with one shared configuration across all four series at the series's native horizon. A natural extension is to fit one LightGBM per horizon ($H = 1, 3, 10, 25$) on the full panel rather than per series, sharing learning across series within a horizon. The per-horizon target std grows by roughly $4.5\times$ from H1 to H25, so per-horizon fitting may produce better-calibrated intervals than per-series fitting alone.
  \item \textbf{Feature-aware deep models.} The deep section in \cref{sec:deep} keeps inputs univariate (lagged target only). A natural extension is to concatenate the lagged top-$k$ features from \cref{sec:eda:feature_audit} into the input window. The expectation is that the volatile and high-weight series gain modestly; the near-zero-signal series do not gain at all. This validates that the deep section's univariate result is not leaving feature-driven lift on the table.
  \item \textbf{Hyperparameter sweep on the LightGBM quantile model.} The current $400$-tree, depth-$5$, $\mathtt{lr}{=}0.04$ configuration is moderate-capacity rather than tuned. A small Optuna sweep ($20$--$50$ trials per quantile) on the calibration slice of one fold could improve point-forecast skill on the highest-weight series, where the current $0.865$ likely has further headroom.
  \item \textbf{Ensembles across the three modelling families.} The classical, deep, and LightGBM-quantile families all produce meaningful but non-overlapping signal: classical works through likelihood, deep through nonlinear sequence learning, LightGBM through gradient-boosted feature splits. A weighted average or ridge-stacking ensemble fit on the meta slice $\mathtt{ts\_index} \in (2880, 3240]$ and scored on the canonical holdout could combine the three families without leakage. The meta slice was reserved untouched precisely so this experiment can be run.
  \item \textbf{Multi-horizon coupling.} The four horizons in the panel are not independent: the targets at $H \in \{3, 10, 25\}$ are well approximated by rolling cumulative sums of the $H = 1$ process (\cref{sec:cumulation}). A multi-output model that predicts all four horizons jointly under a soft cumulation constraint $\hat y^{H_k}_t \approx \sum_i \hat y^{H_1}_{t+i}$ could exploit this latent structure, which the per-series and per-horizon models in this project ignore.
  \item \textbf{Cross-series transfer.} The 86 anonymised features and the cross-series correlations in \cref{sec:eda:cross_series} suggest there is shared structure across sub-codes within a code. A multi-task model that shares a learned representation across sub-codes could exploit this in a way the current per-series classical, per-series deep, and per-series LightGBM-quantile models all do not.
  \item \textbf{Real-money risk management overlay.} The position-sizing rule sketched in \cref{sec:actions:sizing} has not been backtested as a trading strategy. Wrapping the model output in a portfolio-construction layer with hard limits, transaction-cost models, and live coverage monitoring would close the loop from forecast to deployable strategy.
  \item \textbf{Tube Loss as an alternative prediction interval estimator.} The current \cref{sec:quantile} uses LightGBM quantile regression with split conformal post-calibration --- a two-stage pipeline where the model and the coverage guarantee are decoupled. Tube Loss~\cite{anand2024tubeloss} offers a single-stage alternative: a piecewise-linear training objective that jointly optimises a neural network to output lower and upper bounds directly, with an asymptotic coverage guarantee and no sigmoid approximation required. A tunable shift parameter $r \in (0,1)$ allows the interval to be displaced toward denser probability mass, which is especially useful when the conditional distribution of the target is skewed --- as is the case for the most-volatile and highest-weight series in this panel. A width-penalty parameter $\delta$ controls the coverage--width trade-off within a single optimisation. Applied to the LSTM or GRU architectures already evaluated in \cref{sec:deep}, Tube Loss would provide a direct comparison against the current conformal pipeline, testing whether an end-to-end differentiable interval objective can match or exceed the calibrated LightGBM intervals reported in \cref{tab:coverage_summary}.
\end{enumerate}
